%% kapa: exam=1 round=31 year=2020 period=1 form=A subject=civil_law subject_ko=민법 seq=25 number=25
%% source: Q-Net 감정평가사 기출문제 (cst00309) artlSeq=5207219 file=extracted/1008/5207219/A형/A형/2020년도 제31회 감정평가사 1차 시험_1교시_A형.pdf
%% conversion: kordoc 4.24.51 -> tools/gen1.py (auto); review=proofread-sonnet; flags=none
\begin{question}{25}
甲 소유의 X토지에 乙 명의로 소유권이전청구권을 보전하기 위한 가등기를 한 경우에 관한 설명으로 옳은 것은? (다툼이 있으면 판례에 따름)
\begin{choices}
  \item 乙은 부기등기의 형식으로는 가등기된 소유권이전청구권을 양도하지 못한다.
  \item 가등기가 있으면 乙이 甲에게 소유권이전을 청구할 법률관계가 있다고 추정된다.
  \item 乙이 가등기에 기하여 본등기를 하면 乙은 가등기한 때부터 X토지의 소유권을 취득한다.
  \item 가등기 후에 甲이 그의 채권자 丙에게 저당권을 설정한 경우, 가등기에 기하여 본등기를 마친 乙은 丙에 대하여 물상보증인의 지위를 가진다.
  \item 乙이 별도의 원인으로 X토지의 소유권을 취득한 때에는, 특별한 사정이 없으면 가등기로 보전된 소유권이전청구권은 소멸하지 않는다.
\end{choices}
\end{question}
\answer{5}
